% !PRINTED 278
\weberpage{278}
%<<CONT
\[
(-1)^k x^k,
\]

und wenn der absolute Werth von $x$ hinlänglich gross ist, so wird das Vorzeichen dieses Gliedes über das Vorzeichen von $L_k(x)$ entscheiden. Nehmen wir also für $\alpha$ einen genügend grossen negativen, für $\beta$ einen genügend grossen positiven Werth, so finden in (3) für $x=\alpha$ lauter Zeichenwechsel, für $x=\beta$ lauter Zeichenfolgen statt. Es werden also $n$ Wurzeln zwischen $\alpha$ und $\beta$ liegen, und daraus folgt der Satz:

II.\ \gesperrt{Die Gleichung $L_n(x)=0$ hat lauter reelle Wurzeln.}

Diese beiden Sätze sind aber nur von beschränkter Anwendbarkeit, so lange wir uns nicht von der Voraussetzung frei machen können, dass in der Kette der $L_k$ nicht zwei auf einander folgende Glieder zugleich verschwinden sollen. Von dieser Beschränkung können wir den Satz aber durch folgende einfache Ueberlegung befreien.

Nehmen wir an, für irgend einen Werth von $x$ verschwinde $L_k$, aber nicht $L_{k-1}$; wir können dann Grössen $a_{k+1,i}$, die in $L_k$ nicht vorkommen, so bestimmen, dass $L_{k+1}$ für diesen Werth von $x$ nicht verschwindet; denn es kann $L_{k+1}$ nicht identisch für alle $a_{k+1,i}$ verschwinden, weil es das Glied

\[
-a_{k+1,k}^2 L_{k-1}
\]

enthält [\S.\,23, (12)].

Hiernach können wir, wenn in der Kette

\weq{5}{L_n, \; -L_{n-1}, \; L_{n-2} \cdots \pm L_1, \; \mp 1}

einige auf einander folgende Glieder für irgend einen Werth von $x$ verschwinden, durch Abänderung der $a_{i,k}$ eine andere Reihe

\weq{6}{L_n', \; -L_{n-1}', \; L_{n-2}' \cdots \pm L_1', \; \mp 1}

ableiten, in der keine zwei auf einander folgenden Glieder für irgend einen Werth $x$ (des Intervalles $\alpha \ldots \beta$) verschwinden.

Zugleich können die $a_{i,k}'$, von denen die $L'$ abhängen, so angenommen werden, dass sie sich von den $a_{i,k}$ \gesperrt{um weniger als eine beliebig gegebene Grösse $\omega$ unterscheiden}.

Wenn nun die Zahlen $\alpha, \beta$ so angenommen sind, dass in der Reihe (5) kein Glied für $x=\alpha$ oder $x=\beta$ verschwindet, so können wir $\omega$ so klein annehmen, dass entsprechende Glieder von (5) und (6) dasselbe Vorzeichen haben. Durch unseren Satz ist aber die Anzahl der Wurzeln von $L_n'=0$ zwischen $\alpha$ und $\beta$ durch die Zeichen der Reihe (6) bestimmt.
